\begin{center}
    \section*{Santrumpų ir terminų sąrašas}
\end{center}
\addcontentsline{toc}{section}
{Santrumpų ir terminų sąrašas}

\textbf{Santrumpos:} \\
\textbf{pvz.} -- pavyzdžiui. \\
\textbf{angl.} -- angliškai. \\
\textbf{CI/CD} -- Nuolatinė integracija / Nuolatinis pristatymas (angl. \textit{Continuous Integration / Continuous Delivery}). \\
\textbf{HTTP/HTTPS} -- HyperText Transfer Protocol [Secure] (\textit{liet.} hiperteksto perdavimo protokolas). \\
% \textbf{XML} -- eXtensible Markup Language (\textit{liet.} išplečiama žymėjimo kalba). \\
% \textbf{JSON} -- JavaScript Object Notation (\textit{liet.} JavaScript objektų žymėjimas). \\
% \textbf{API} -- Application Programming Interface (\textit{liet.} programų programavimo sąsaja). \\
% \textbf{UML} -- Unified Modeling Language. \\
% \textbf{WSL} -- Windows Subsystem for Linux (\textit{liet.} Linux  posistemė Windows sistemoje). \\
% \textbf{SSL/TLS} -- Secure Sockets Layer / Transport Layer Security. \\
% \textbf{WebDAV} -- Web-based Distributed Authoring and Versioning. \\
% \textbf{E2E} -- End-to-End testing.
% \textbf{WSL} -- Windows Subsystem for Linux (\texit{liet.} Linux  posistemė Windows sistemoje). \\




\textbf{Terminai:} \\
\textbf{Pažymėjimas} (angl. \textit{annotation} arba \textit{highlight}) -- knygos teksto ištrauka, kurią naudotojas nori išsaugoti ir grįžti vėliau. \\
\textbf{Pažymėjimo pastaba} (angl. \textit{note}) -- naudotojo komentaras apie konkrečią teksto ištrauką; komentaras yra saugomas kartu su pažymėjimu. \\
\textbf{Uždara ekosistema} (angl. \textit{vendor-lock-in}) -- sistema, kurioje vartotojai yra apriboti tam tikroje platformoje, negalėdami lengvai perkelti savo duomenų ar naudotis kitų gamintojų paslaugomis. \\

\newpage